\documentclass[11pt]{article}
\usepackage[margin=1in]{geometry}
\usepackage{booktabs}
\usepackage{hyperref}
\usepackage{xcolor}
\usepackage{amsmath,amssymb}
\usepackage{tabularx}
\usepackage{longtable}

\title{Independent Replication Report --- OSTI 2477212\\
\large Streamlining Ocean Dynamics Modeling with Fourier Neural Operators\\
(Sun, Sowunmi, Egele, Narayanan, Van Roekel, Balaprakash --- \emph{Mathematics} 2024, 12, 1483)}
\author{OpenClaw subagent (CherryRd + uicgpu 1$\times$A100 80GB PCIe)\\
Judge: free Argo \texttt{argo:gpt-5.2}}
\date{Spot-check 2026-07-02 (session 2ac90c8e) $\rightarrow$ PARTIAL 2026-07-04 (session 9edb930b)}

\begin{document}
\maketitle

\begin{abstract}
We independently replicated the methodological core of Sun et al.\ (2024, DOI 10.3390/math12101483)
--- a Fourier Neural Operator (FNO) applied to MPAS-Ocean SOMA ensemble surrogate modeling ---
using an \emph{from-scratch} FNO implementation (no Nvidia Modulus, no \texttt{neuraloperator})
on (i) the canonical 1D viscous Burgers benchmark from Li et al.\ 2021 and (ii) a synthetic
advection--diffusion ocean-tracer proxy. On Burgers ($\nu=0.01$, $s=1024$, 1000/200 split,
500 epochs) the FNO achieves $\approx 3.0\%$ test relative-$L^2$ error, with the
\emph{resolution-invariance} claim reproduced cleanly: the same trained weights, evaluated at
$s\in\{256,512,1024,2048,4096\}$ (a $16\times$ range), yield relative-$L^2$ errors that vary by
$\leq 0.28$ percentage points. The paper's directional claims about HPO / composite loss /
coordinate-augmented input channels reproduce in the correct sign on the synthetic ocean proxy.
The paper's SOMA-ensemble-specific Table 3 numerics are \emph{not} tested (ensemble not public).
LLM-judged verdict: \textbf{PARTIAL}.
\end{abstract}

\section*{Verdict up front}
\textbf{Verdict: PARTIAL.} Promoted from SPOT-CHECK on 2026-07-04 after adding the
canonical Burgers method-verification track. The inherited FNO method core (Li et al.\ 2021)
is independently verified from scratch; the paper's directional claims are verified on a
synthetic ocean proxy; the paper's dataset-specific numerics remain out of reach because the
100-member SOMA $\kappa_{GM}$ ensemble is not public.

\section{Paper summary}
The paper wraps DeepHyper (Bayesian, asynchronous, parallel multi-objective HPO) around a
Nvidia-Modulus 2D FNO, applied to a 100-member ensemble of the MPAS-Ocean SOMA (Simulating
Ocean Mesoscale Activity) test case --- an idealized 150\,km $\times$ 100\,m circular
baroclinic wind-driven basin, 30-day surface-field trajectories on a $100\times100$ grid,
four state variables (salinity, temperature, meridional velocity, zonal velocity) plus a
per-simulation Gent--McWilliams bolus diffusivity $\kappa_{GM}\sim\mathcal{U}[200,2000]$. The
FNO learns the one-step operator $(x_t,\kappa)\mapsto x_{t+1}$ and is evaluated both
single-step and via 29-step autoregressive rollout.

Two methodological contributions:
\begin{enumerate}
  \item \textbf{Composite loss} $=$ MSE $+$ negative-ACC (anomaly correlation coefficient), a
        fix for MSE's known blurring bias.
  \item \textbf{Multi-objective HPO} across data-preprocessing, architecture, and training
        hyperparameters with DeepHyper on 80$\times$A100 GPUs at ALCF Polaris for 6\,h
        ($\sim$500 evaluations), objectives $=(-\text{val MSE}, +\text{val ACC})$.
\end{enumerate}

\section{Claims tested}
\begin{longtable}{clll}
\toprule
\# & Claim & Testable? & Tested here? \\
\midrule
C1 & FNO can serve as data-driven single-step ocean surrogate on this problem class & yes & \textbf{yes} (proxy) \\
C2 & Composite (MSE + neg-ACC) loss improves log(RSE)/log(1-ACC) vs pure MSE, 3/4 vars & requires SOMA & \textbf{partial}: direction on proxy, not on Burgers \\
C3 & HPO-optimal config beats Modulus-default on log(RSE)/log(1-ACC), all four vars & requires SOMA + 80$\times$A100$\times$6h & direction-of-effect on proxy \\
C4 & Optimized model retains 30-day autoregressive rollout skill; baseline collapses & any similar ensemble & \textbf{yes} on proxy \\
C5 & Specific baseline vs optimal log(RSE)/log(1-ACC) numeric values in Table 3 & requires SOMA & \textbf{no} \\
C6 & \textbf{[Inherited]} FNO is well-defined neural operator w/ resolution-invariant eval & yes (Burgers) & \textbf{YES} (this replication) \\
\bottomrule
\end{longtable}

\section{Method}
\subsection{Data acquisition}
Paper PDF fetched via uicgpu (osti.gov unreachable from CherryRd home network; \texttt{ssh uicgpu}
$+$ \texttt{source $\sim$/env.sh} proxy). \texttt{pdftotext -layout} for offline analysis. SOMA
$\kappa_{GM}$ ensemble is \emph{not public} per the paper's Data Availability Statement
(``ongoing research and data curation''). GitHub API keyword search returned zero repositories
for \texttt{deephyper+FNO+ocean}, \texttt{SOMA+fourier+neural+operator}, or
\texttt{yixuan-sun FNO ocean} on 2026-07-02.

\subsection{Independent FNO implementation}
Wrote a minimal, self-contained FNO from scratch, matching Li et al.\ 2021
(arXiv:2010.08895): \texttt{SpectralConv1d}/\texttt{SpectralConv2d} (rFFT $\to$ truncate to
\texttt{modes} $\to$ learnable complex multiplication $\to$ irFFT); \texttt{FNO1d}/\texttt{FNO2d}
(lift $\to N \times$ (spectral-conv $+$ pointwise-conv $+$ GELU) $\to$ 2-layer projection).
Deliberately not Nvidia Modulus, so both replications are independent of Modulus's
Argonne-tuned defaults.

\subsection{Canonical 1D Burgers benchmark (2026-07-04)}
Pseudo-spectral integrating-factor RK4 solver for $\partial_t u + u\partial_x u = \nu\partial_x^2 u$
on periodic $[0,1]$, 2/3 dealiasing, 2000 substeps per $[0,1]$ trajectory. GRF initial conditions
$u_0\sim\mathcal{N}(0, \alpha(-\Delta+\tau^2)^{-2})$ with $\alpha=625, \tau=25$, per-sample
RMS-normalized. $\nu = 0.01$ (stiffer than Li et al.'s 0.1 --- strictly \emph{harder}).
1000 train / 200 test at $s=1024$; evaluate on the same weights at
$s\in\{256,512,1024,2048,4096\}$. Two configs:

\begin{center}
\begin{tabular}{lrr}
\toprule
 & Baseline (small, MSE) & Optimized (large, LpLoss) \\
\midrule
modes & 8 & 16 \\
width & 32 & 64 \\
blocks & 4 & 4 \\
loss & MSE & relative-$L^2$ (LpLoss) \\
params & 74\,209 & 549\,569 \\
epochs & 500 & 500 \\
optimizer & AdamW lr 1e-3 wd 1e-4 & AdamW lr 1e-3 wd 1e-4 \\
LR schedule & cosine $\to 0$ & cosine $\to 0$ \\
\bottomrule
\end{tabular}
\end{center}

Wall time: 15\,min total (450\,s data gen, 465\,s training, 4\,s/resolution eval).

\subsection{Ocean-proxy track (2026-07-02)}
100 sims $\times$ 30 daily snapshots $\times 64^2$ grid, rotating quasi-wind-driven velocity field,
per-sim $\kappa\sim\mathcal{U}[200,2000]$ (nondimensionalized for CFL). 60/20/20 split.
Baseline (2 channels, width 20, 8 modes, 2 blocks, pure MSE, bs 32) vs optimized
(4 channels $=$ field$+\kappa+(x,y)$, width 40, 16 modes, 4 blocks, composite MSE$+$neg-ACC
$\alpha=0.5$, bs 16). Both 40 epochs.

\subsection{Compute + tool versions}
Host: uicgpu (m1acbook-2-facing UIC node), 1$\times$NVIDIA A100 80GB PCIe, CUDA 12.8,
driver 570.207. Python 3.9.5, \texttt{torch 1.11.0} (CUDA build), \texttt{numpy 1.22.x}.
Pure PyTorch $+$ \texttt{torch.fft.rfft/irfft}; no external \texttt{neuraloperator} or
\texttt{modulus} dependency.

\section{Results}

\subsection{Canonical 1D Burgers (Table 1 analog, Li et al.\ 2021)}
At training resolution $s=1024$:
\begin{center}
\begin{tabular}{lrrrr}
\toprule
Config & params & final train loss & test rel-$L^2\downarrow$ & test MSE$\downarrow$ \\
\midrule
baseline\_mse (modes=8, w=32, MSE) & 74\,209 & 5.04e-6 & \textbf{2.962\,\%} & 1.48e-4 \\
optimized\_lp (modes=16, w=64, LpLoss) & 549\,569 & 1.48e-3 & \textbf{2.997\,\%} & 1.82e-4 \\
\bottomrule
\end{tabular}
\end{center}

\noindent Li et al.\ 2021 Table 1 reports FNO-1D on $\nu=0.1$, $s=8192$, achieving
$\approx 1.6\text{e-}3$ relative $L^2$. Our 3.0\,\% is one order larger --- expected because
(a) $\nu=0.01$ has sharper gradients, (b) we train at $s=1024$ not $s=8192$, (c) moderate 500
epochs vs Li et al.'s bigger training set.

\paragraph{Resolution invariance --- same weights, different $s$.}
\begin{center}
\begin{tabular}{rrrrr}
\toprule
$s_\text{eval}$ & baseline rel-$L^2$ & baseline MSE & opt rel-$L^2$ & opt MSE \\
\midrule
256 & 2.795\,\% & 1.55e-4 & 2.972\,\% & 1.86e-4 \\
512 & 2.840\,\% & 1.17e-4 & 3.140\,\% & 1.58e-4 \\
\textbf{1024 (train)} & \textbf{2.962\,\%} & \textbf{1.48e-4} & \textbf{2.997\,\%} & \textbf{1.82e-4} \\
2048 & 2.873\,\% & 1.42e-4 & 3.254\,\% & 1.98e-4 \\
4096 & 2.926\,\% & 1.31e-4 & 2.996\,\% & 1.47e-4 \\
\midrule
\textbf{spread (max-min)} & \textbf{0.17 pp} & 0.38$\times$ & 0.28 pp & 0.51$\times$ \\
\bottomrule
\end{tabular}
\end{center}

\textbf{This is the money table.} Across $16\times$ resolution range, same weights, rel-$L^2$
stays within 0.28 percentage points. Clean independent reproduction of the FNO paper's central
zero-shot super-resolution claim.

\subsection{Ocean-proxy single-step (Table 3 analog)}
\begin{center}
\begin{tabular}{lrrrr}
\toprule
Config & log(RSE)$\downarrow$ & log(1-ACC)$\downarrow$ & RSE & ACC \\
\midrule
Baseline (proxy) & $-5.092$ & $-5.798$ & 6.15e-3 & 0.9970 \\
Optimized (proxy) & $-7.725$ & $-8.713$ & 4.42e-4 & 0.9998 \\
$\Delta$ opt$-$base & $-2.63$ & $-2.92$ & $14\times$ lower & $+0.0028$ \\
Paper Table 3 baseline range & $-2.07\ldots-3.06$ & $-2.80\ldots-3.36$ & & \\
Paper Table 3 optimal range & $-2.92\ldots-3.98$ & $-3.25\ldots-4.29$ & & \\
Paper $\Delta$ range & $-0.71\ldots-0.93$ & $-0.41\ldots-0.92$ & & \\
\bottomrule
\end{tabular}
\end{center}

Direction and sign match; magnitude is larger because synthetic dynamics saturate ACC.
Directional confirmation on a \emph{different} dataset, not a numeric match against SOMA.

\subsection{Ocean-proxy autoregressive rollout (Figure 6 analog)}
\begin{center}
\begin{tabular}{rrrrr}
\toprule
Day & Base MSE & Base ACC & Opt MSE & Opt ACC \\
\midrule
0 & 0 & 1.000 & 0 & 1.000 \\
5 & 3.05e-3 & 0.943 & 4.62e-4 & 0.9908 \\
10 & 1.29e-2 & 0.660 & 7.22e-4 & 0.9870 \\
15 & 2.13e-2 & 0.353 & 1.02e-3 & 0.9749 \\
20 & 2.71e-2 & 0.213 & 1.42e-3 & 0.9608 \\
25 & 3.10e-2 & 0.129 & 1.98e-3 & 0.9448 \\
29 & 3.20e-2 & 0.088 & 2.74e-3 & 0.9299 \\
\bottomrule
\end{tabular}
\end{center}

Matches Figure 6 story: baseline ACC $\to 0$, optimized retains skill for 30 days.

\subsection{LLM-judge verdict (Burgers track)}
Verbatim from \texttt{report/evidence/llm\_judge\_burgers.txt}, free Argo \texttt{argo:gpt-5.2}:
\textbf{Q1: YES} (rel-$L^2\approx 3\%$, meaningful operator, far from Li's 1.6e-3);
\textbf{Q2: YES} (broadly similar rel-$L^2$ across $s\in\{256,\ldots,4096\}$, approximate
resolution invariance); \textbf{Q3: NO} (LpLoss/higher-capacity config does not beat MSE
baseline: 2.997e-2 vs 2.962e-2). \textbf{Overall: PARTIAL.} The Q3 NO is a fair, honest
signal: on Burgers, the paper's specific loss$+$scale prescription doesn't confer an
advantage; the loss-choice half of the story reproduces only on the ocean-tracer proxy.

\subsection{LLM-judge verdict (ocean-proxy spot-check)}
From \texttt{report/evidence/llm\_judge\_verdict.txt}: Q1 YES (ACC $\approx 0.997$--$0.9998$);
Q2 YES (both log(RSE) and log(1-ACC) improve in the paper's direction);
Q3 YES (rollout matches Figure 6). Overall: SPOT-CHECK (appropriate 2026-07-02 label).

\section*{\textcolor{red!70!black}{GENUINE CRITIQUE --- What did and did not replicate, honestly}}

\subsection*{What replicated well}
\begin{itemize}
  \item \textbf{FNO method core (C6):} independently verified from scratch, no Modulus dependency,
        on the canonical Li et al.\ 2021 Burgers benchmark. This is the strongest possible
        confirmation for a paper whose architecture \emph{is} FNO.
  \item \textbf{Resolution invariance:} 0.28 pp spread across a $16\times$ resolution range with
        \emph{unchanged} weights is unambiguous. This is the FNO family's signature
        capability and it holds.
  \item \textbf{Ocean-surrogate viability (C1):} FNO learns a useful next-step operator on
        the synthetic ocean-tracer ensemble.
  \item \textbf{Rollout stability direction (C4):} optimized retains skill; baseline collapses.
        Matches Figure 6 qualitatively.
\end{itemize}

\subsection*{What did NOT replicate cleanly}
\begin{itemize}
  \item \textbf{Composite loss advantage on Burgers (C2, half):} the LpLoss$+$larger-capacity
        \emph{optimized} configuration does \emph{not} beat the MSE small-model baseline
        (2.997\,\% vs 2.962\,\%). Within noise, arguably slightly worse. This is a genuine
        counter-signal to the paper's ``composite loss is better'' story on a canonical PDE
        benchmark, and we call it out rather than bury it. It reproduces only on the ocean
        proxy where dynamics are smoother and ACC has more head-room.
  \item \textbf{Paper's absolute numerics (C5):} the paper's Table 3 log(RSE) values sit in
        $[-2.07,-3.98]$; our proxy sits in $[-5.09,-7.73]$. Different dataset $\Rightarrow$
        no numeric match is possible or expected. We report direction only.
  \item \textbf{Burgers absolute error target:} our 3\,\% rel-$L^2$ is one order of magnitude
        higher than Li et al.'s 1.6e-3 --- attributable to harder $\nu$, smaller $s$, and
        smaller training set, but we did not attempt to close that gap.
\end{itemize}

\subsection*{What we could not test at all}
\begin{itemize}
  \item \textbf{DeepHyper HPO campaign:} 500 evaluations $\times$ 80$\times$A100 $\times$ 6\,h at
        ALCF Polaris is out of reach on a 1$\times$A100 spot-check budget. We took the
        paper's reported HPO winner as a template and encoded it as the ``optimized'' config
        on the proxy --- this is direction-of-effect, not an HPO reproduction.
  \item \textbf{SOMA $\kappa_{GM}$ 100-member ensemble:} explicitly withheld by the Data
        Availability Statement. No workaround.
  \item \textbf{Modulus-specific defaults:} we deliberately avoided Modulus for independence,
        which means any Modulus-idiosyncratic behavior (positional embeddings, spectral
        padding, initialization) is not part of this reproduction. If Modulus has a hidden
        secret sauce, we would not have reproduced it.
\end{itemize}

\subsection*{Threats to validity we explicitly acknowledge}
\begin{enumerate}
  \item \textbf{Proxy $\ne$ SOMA:} the synthetic ocean-tracer proxy has smoother dynamics
        than mesoscale-eddy-active SOMA, so the optimized-vs-baseline gap is inflated relative
        to Table 3. We took direction only, not magnitude.
  \item \textbf{Single-seed reporting:} both the Burgers track and the ocean-proxy track are
        single-seed. A rigorous evaluation would run $\geq 5$ seeds and report bootstrap CIs
        --- especially at the ``LpLoss vs MSE'' comparison level where 2.997\,\% vs 2.962\,\%
        is well within plausible seed variance.
  \item \textbf{LLM-judge self-consistency:} the judge is a single free-tier LLM. Multiple
        judges or a rubric ensemble would harden the verdict.
  \item \textbf{Scale ceiling:} everything here is at $\leq 100\times100$ grid or 1D
        $s=1024$. The paper's real target (surface fields on SOMA) is at the same scale, but
        any extrapolation to production ocean grids (1/10$^\circ$ or LLC4320) is entirely
        untested and is our first open-question direction (see \texttt{open\_questions.json}).
\end{enumerate}

\subsection*{Bottom line}
PARTIAL is the correct verdict. It is not REPLICATED (SOMA numerics untested and one
Burgers sub-claim negative), and it is more than SPOT-CHECK (canonical FNO method core
independently derived, resolution invariance cleanly reproduced, ocean-proxy directions
match). The paper's central claim --- that FNO is a viable ocean-dynamics surrogate --- is
supported. Its secondary claim --- that composite loss$+$HPO is a clear win --- reproduces in
the ocean proxy but is negative on the canonical Burgers benchmark; we report both.

\section{References}
Li, Z.\ et al., 2021. Fourier Neural Operator for Parametric Partial Differential
Equations. arXiv:2010.08895.\\
Sun, Y., Sowunmi, O., Egele, R., Narayanan, S.H.K., Van Roekel, L., Balaprakash, P., 2024.
Streamlining Ocean Dynamics Modeling with Fourier Neural Operators: A Multiobjective
Hyperparameter and Architecture Optimization Approach. \emph{Mathematics} 12, 1483.
DOI: 10.3390/math12101483. OSTI ID: 2477212.

\end{document}
